%% Chapter 14: Approximation Algorithms II
%% Status: drafted; figures are still placeholders.
%%
%% Written against slides 30-74 of "material/slides/Approximation Algorithms 2026.pdf"
%% (Johan van Rooij).  The pool exercise on Steiner tree is worked in the last
%% section before the exercises; it is not in the deck, but it is in the pool,
%% and both halves of it are applications of techniques that are.
%%
%% Material deliberately NOT here because it is not in the deck: the
%% primal-dual method, randomised rounding, APX-completeness (that is
%% chapter 16).

\chapter{Approximation Algorithms II}
\label{ch:approx-ii}

\Cref{ch:approx-i} had one quality of approximation: a constant factor. That is
the middle of a scale. Some problems can be solved to within an additive
constant, which is far better; some can be approximated to within
$1+\varepsilon$ for any $\varepsilon$, which is nearly as good; some resist any
constant and admit only a ratio growing with the input. This chapter lays out
the scale, gives an algorithm at each level of it, and --- just as important ---
gives the techniques for proving that a problem does \emph{not} sit higher on
the scale than it does.

\begin{goals}
  \poolgoal{name the five qualities of polynomial-time approximation and
        give an example of each}
  \poolgoal{give a problem approximable within an absolute constant
        difference, and show that maximum independent set is not one unless
        $\mathrm{P} = \mathrm{NP}$}
  \poolgoal{define PTAS and FPTAS, and give an FPTAS for \lang{Knapsack}:
        the dynamic program, the scaling, the choice of the scaling constant,
        and the proofs of ratio and running time}
  \poolgoal{use the gap technique to show that \lang{Knapsack} admits no
        absolute constant difference approximation unless
        $\mathrm{P} = \mathrm{NP}$}
  \poolgoal{give the PTAS for maximum independent set on planar graphs,
        and prove its ratio and running time}
  \poolgoal{explain the polynomial-bound technique, and use it to show
        that maximum independent set admits no FPTAS unless
        $\mathrm{P} = \mathrm{NP}$}
  \poolgoal{give the greedy $(\ln n + 1)$-approximation for
        \lang{Set Cover} and prove its ratio}
  \poolgoal{give the self-reduction from $G$ to $G \times G$ and use it to
        show that maximum independent set is not in APX unless it has a PTAS}
  \poolgoal{prove that the shortest-path-metric MST heuristic is a
        $2$-approximation for \lang{Steiner Tree}}
  \poolgoal{prove that there is no FPTAS for \lang{Steiner Tree} unless
        $\mathrm{P} = \mathrm{NP}$}
\end{goals}

%==============================================================================
\section{Five qualities of approximation}
\label{sec:approx-qualities}
%==============================================================================

All five of the following ask for a polynomial-time algorithm. They differ in
what is promised about the answer, and they are listed from best to worst.

\begin{enumerate}[leftmargin=*,label=\arabic*.]
  \item \textbf{Absolute constant difference.}
        $|\OPT - \mathrm{ALG}| \le c$ for a constant $c$. The error does not
        grow with the instance at all.
  \item \textbf{FPTAS}, a fully polynomial-time approximation scheme. For every
        $\varepsilon > 0$ the algorithm achieves ratio $1+\varepsilon$, and its
        running time is polynomial in the input size \emph{and} in
        $1/\varepsilon$.
  \item \textbf{PTAS}, a polynomial-time approximation scheme. For every fixed
        $\varepsilon > 0$ the algorithm achieves ratio $1+\varepsilon$ in
        polynomial time --- but the exponent of the polynomial may depend on
        $\varepsilon$, so the running time may be $n^{1/\varepsilon}$.
  \item \textbf{APX.} Ratio $c$ for some constant $c$, and no better.
  \item \textbf{$f(n)$-APX.} Ratio $f(n)$, a function of the input size only.
\end{enumerate}

\begin{definition}[PTAS, FPTAS]\label{def:ptas-fptas}
An approximation \emph{scheme} for an optimisation problem is an algorithm
taking two inputs, an instance $X$ and a real $\varepsilon > 0$, that returns a
solution of value within a factor $1+\varepsilon$ of optimal. It is a
\emph{polynomial-time approximation scheme} (PTAS) if for every fixed
$\varepsilon > 0$ it runs in time polynomial in $|X|$, and a \emph{fully
polynomial-time approximation scheme} (FPTAS) if it runs in time polynomial in
$|X|$ and $1/\varepsilon$ together.
\end{definition}

The single word \emph{fully} is the whole difference, and it matters. A PTAS is
allowed a running time of $n^{1/\varepsilon}$, which for $\varepsilon = 0.01$ is
$n^{100}$ and useless; an FPTAS must stay polynomial in $1/\varepsilon$, for
instance $n^3/\varepsilon$, which is usable. Every FPTAS is a PTAS, and every
problem with a PTAS is in \textrm{APX} --- take $\varepsilon = 1$ --- so the
list above really is a chain of inclusions.

Each level of the chain has a matching negative technique, and the negative
results are what make the classification worth having:
\begin{itemize}[leftmargin=*,nosep]
  \item the \emph{gap technique} of \cref{sec:gap-technique} separates
        \textrm{APX} from $f(n)$-APX, and also rules out absolute constant
        differences (\cref{sec:knapsack-gap});
  \item the \emph{polynomial-bound technique} (\cref{sec:polynomial-bound})
        separates FPTAS from PTAS;
  \item \emph{APX-hardness} (\cref{ch:hardness}) separates PTAS from
        \textrm{APX}.
\end{itemize}

%==============================================================================
\section{Absolute constant difference}
\label{sec:absolute-difference}
%==============================================================================

The best guarantee short of optimality is an error that does not grow at all.
Two examples, in both of which a hard theorem does the work.

\begin{example}[Colouring a planar graph]\label{ex:planar-colouring}
By the four colour theorem every planar graph is $4$-colourable, and testing
$1$- and $2$-colourability is easy. So: test whether $G$ is $1$-colourable, then
whether it is $2$-colourable, and otherwise output $4$. The answer is wrong only
when $G$ is $3$-chromatic, and then by exactly one. This is an approximation
with absolute difference $1$ --- and it is worth noticing that deciding whether
a planar graph is $3$-colourable is \textrm{NP}-complete, so the one case the
algorithm gets wrong is precisely the hard one.
\end{example}

\begin{example}[Edge colouring]\label{ex:vizing}
Vizing's theorem says that every graph $G$ can be edge coloured with at most
$\Delta(G)+1$ colours, where $\Delta(G)$ is the maximum degree, and clearly at
least $\Delta(G)$ colours are needed. Returning $\Delta(G)+1$ is therefore an
approximation of the chromatic index with absolute difference $1$.
\end{example}

Such examples are rare, and there is a reason. For most problems an additive
constant can be amplified away.

\begin{theorem}\label{thm:no-additive-is}
Unless $\mathrm{P} = \mathrm{NP}$, there is no polynomial-time algorithm that
approximates maximum independent set within an absolute difference $c$, for any
integer constant $c > 0$.
\end{theorem}

\begin{proof}
Suppose such an algorithm exists. Given a graph $G$ and an integer $k$, build
$G'$ consisting of $c+1$ disjoint copies of $G$ and run the algorithm on $G'$.

An independent set of $G'$ is a union of independent sets of the copies, so
$\OPT(G') = (c+1)\OPT(G)$.

If $G$ has an independent set of size $k$, then $\OPT(G') \ge (c+1)k$ and the
algorithm returns a set of size at least $(c+1)k - c$. That set is spread over
$c+1$ copies, so some copy contributes at least
$\lceil ((c+1)k-c)/(c+1) \rceil = k$ vertices, giving an independent set of size
$k$ in $G$.

If $G$ has no independent set of size $k$, then $\OPT(G) \le k-1$, so the
algorithm returns at most $\OPT(G') = (c+1)(k-1) < (c+1)k - c$.

The two cases are distinguished by the size returned, so we have decided an
\textrm{NP}-complete problem in polynomial time.
\end{proof}

\begin{intuition}
Disjoint copies multiply the optimum and leave the additive error alone. Any
problem whose instances can be combined so that the optima add --- and most can
--- kills additive approximation the same way. This is why the two examples
above come from theorems of the form ``the answer is always one of two
consecutive numbers'': there the additive constant is not a guarantee that was
earned by an algorithm, it is a fact about the problem.
\end{intuition}

%==============================================================================
\section{An FPTAS for knapsack}
\label{sec:knapsack-fptas}
%==============================================================================

The example at level~2 of the scale is knapsack, and the route to it runs
through an algorithm that is not polynomial at all.

\begin{probdef}{Knapsack}
  \instance $n$ items, item $i$ having a positive integer weight $w(i)$ and a
  positive integer value $v(i)$, and an integer capacity $W$.
  \question A set of items of total weight at most $W$ whose total value is as
  large as possible.
\end{probdef}

\subsection{A pseudo-polynomial algorithm}

Let $P = \max_i v(i)$, so that no solution has value more than $nP$. Tabulate
\[ M(i,Q) = \text{the minimum total weight of a subset of items } 1,\dots,i
   \text{ of total value exactly } Q, \]
for $0 \le i \le n$ and $0 \le Q \le nP$, with $M(i,Q) = \infty$ when no such
subset exists. The recurrence is
\[ M(0,0) = 0, \qquad M(0,Q) = \infty \ (Q > 0), \qquad
   M(i+1,Q) = \min\bigl\{ M(i,Q),\; M(i,Q-v(i+1)) + w(i+1) \bigr\}, \]
where the second term is omitted if $Q < v(i+1)$. The answer is the largest $Q$
with $M(n,Q) \le W$. The table has $\bigoh{n^2 P}$ entries and each costs
constant time, so the algorithm runs in $\bigoh{n^2 P}$ time.

\begin{definition}[Pseudo-polynomial]\label{def:pseudo-polynomial}
An algorithm is \emph{pseudo-polynomial} if its running time is polynomial in
the input size and in the largest number appearing in the input.
\end{definition}

That is not a polynomial-time algorithm: numbers are written in binary, so $P$
can be exponential in the input size, and then $\bigoh{n^2P}$ is exponential
too. It is polynomial exactly when $P$ is small --- and the idea of the next
subsection is to \emph{make} $P$ small, by throwing away the low-order bits of
the values, and to pay for the damage in the approximation ratio.

\subsection{Scaling}

\begin{algorithm}[ht]
\caption{FPTAS for \lang{Knapsack}}
\label{alg:knapsack-fptas}
\begin{algorithmic}[1]
  \State discard every item with $w(i) > W$, since no solution uses one
  \State let $c > 0$ be a scaling constant, fixed below
  \State build a new instance with the same weights and values
         $v'(i) = \lfloor v(i)/c \rfloor$
  \State solve the new instance optimally with the dynamic program above
  \State \Return that set of items, as a solution to the original instance
\end{algorithmic}
\end{algorithm}

Because the weights are untouched, whatever the dynamic program returns is
feasible for the original instance. Two questions remain: how much value is
lost, and how long it takes.

\begin{lemma}\label{lem:knapsack-scaling}
\Cref{alg:knapsack-fptas} returns a solution of value at least $\OPT - nc$, and
runs in time $\bigoh{n^2P/c}$.
\end{lemma}

\begin{proof}
For the running time, the largest scaled value is at most $P/c$, so the table of
the dynamic program has $\bigoh{n \cdot nP/c}$ entries.

For the value, let $Y$ be an optimal solution of the original instance, of value
$\OPT$. Each item loses less than one unit of scaled value in the rounding,
since $v(i)/c - \lfloor v(i)/c\rfloor < 1$, and $Y$ has at most $n$ items, so
$Y$ has scaled value more than $\OPT/c - n$. The dynamic program returns a
scaled-optimal solution, so the set it returns has scaled value at least
$\OPT/c - n$ as well. Each unit of scaled value is worth at least $c$ units of
real value, since $v(i) \ge c\lfloor v(i)/c \rfloor$, so the real value of the
returned set is at least $c(\OPT/c - n) = \OPT - nc$.
\end{proof}

\begin{theorem}\label{thm:knapsack-fptas}
With $c = \varepsilon P / (2n)$, \cref{alg:knapsack-fptas} is a
$(1+\varepsilon)$-approximation running in $\bigoh{n^3/\varepsilon}$ time for
every $0 < \varepsilon \le 1$. \lang{Knapsack} therefore has an FPTAS.
\end{theorem}

\begin{proof}
Running time: substituting into \cref{lem:knapsack-scaling},
\[ \bigoh{\frac{n^2 P}{c}} = \bigoh{\frac{n^2P}{\varepsilon P/(2n)}}
   = \bigoh{\frac{n^3}{\varepsilon}}, \]
which is polynomial in $n$ and in $1/\varepsilon$, as an FPTAS requires. Note
that $P$ has cancelled: that is the point of the scaling.

Ratio: since items heavier than $W$ were discarded, every remaining single item
is by itself a feasible solution, so $\OPT \ge P$. By
\cref{lem:knapsack-scaling},
\[ \mathrm{ALG} \;\ge\; \OPT - nc \;=\; \OPT - \frac{\varepsilon P}{2}
   \;\ge\; \OPT - \frac{\varepsilon}{2}\OPT, \]
and therefore
\[ \frac{\OPT}{\mathrm{ALG}} \;\le\; \frac{\OPT}{\OPT(1 - \varepsilon/2)}
   \;=\; \frac{1}{1-\varepsilon/2} \;\le\; 1 + \varepsilon, \]
where the last step uses $\varepsilon \le 1$.
\end{proof}

\begin{intuition}
The scaling constant is chosen by looking at the two bounds and making them
meet. \Cref{lem:knapsack-scaling} says the loss is $nc$ and the time is
$n^2P/c$: raising $c$ buys time and costs accuracy. The loss only matters
relative to $\OPT$, and $\OPT \ge P$, so a loss of $\varepsilon P/2$ is
affordable --- which fixes $nc = \varepsilon P/2$ and hence $c$. Every FPTAS
obtained by rounding has this shape: find the quantity the error can be
measured against, and round away everything far below it.
\end{intuition}

%==============================================================================
\section{Knapsack and the gap technique}
\label{sec:knapsack-gap}
%==============================================================================

\Cref{thm:knapsack-fptas} puts \lang{Knapsack} at level~2 of the scale. Is it at
level~1? No, and the gap technique of \cref{sec:gap-technique} settles it in
four lines.

\begin{theorem}\label{thm:knapsack-no-additive}
Unless $\mathrm{P} = \mathrm{NP}$, no polynomial-time algorithm approximates
\lang{Knapsack} within an absolute difference $c$, for any integer constant
$c \ge 1$.
\end{theorem}

\begin{proof}
Take any instance and multiply every item value by $c+1$, leaving weights and
capacity alone. The feasible sets are unchanged and every solution value is
multiplied by $c+1$; in particular all attainable values are multiples of
$c+1$, so any two distinct values differ by at least $c+1$.

Run the assumed algorithm on the new instance. It returns a solution whose value
is within $c$ of the optimum, and the only attainable value that close to the
optimum is the optimum itself. So the algorithm solves \lang{Knapsack} exactly
in polynomial time, and \lang{Knapsack} is \textrm{NP}-hard.
\end{proof}

The gap here was created by inflation rather than by a reduction from another
problem, but it is the same technique: arrange for the values of interest to be
so far apart that an approximation cannot land between them.

%==============================================================================
\section{A PTAS for independent set on planar graphs}
\label{sec:ptas-planar-is}
%==============================================================================

Now an example at level~3 that is not at level~2. The construction is Baker's,
and it is the model for a large family of approximation schemes in planar and
geometric settings.

\begin{definition}[Levels, $k$-outerplanar]\label{def:k-outerplanar}
Fix a planar embedding of $G$. The vertices on the outer face have \emph{level}
$1$; having removed all vertices of level less than $i$, the vertices then on
the outer face have level $i$. A graph is \emph{$k$-outerplanar} if it has a
planar embedding with at most $k$ levels.
\end{definition}

\begin{figure}[ht]
\centering
\figplaceholder{5cm}{A planar embedding with its vertices grouped into
concentric levels, level $1$ on the outer face, and the same drawing with one
level class deleted, leaving a disjoint union of $k$-outerplanar pieces.}
\caption{Levels of a planar embedding. Deleting every $(k+1)$st level breaks the
graph into $k$-outerplanar components, each of which can be solved exactly.}
\label{fig:planar-levels}
\end{figure}

\begin{lemma}\label{lem:k-outerplanar-is}
For every $k$, maximum independent set restricted to $k$-outerplanar graphs can
be solved optimally in $\bigoh{8^k n}$ time.
\end{lemma}

We do not prove this here. The reason it is true is that a $k$-outerplanar graph
has treewidth at most $3k-1$, and \cref{ch:treewidth} solves independent set on
a graph of treewidth $t$ in $\bigoh{2^t n}$ time.

\begin{algorithm}[ht]
\caption{PTAS for maximum independent set on planar graphs}
\label{alg:planar-is-ptas}
\begin{algorithmic}[1]
  \Require a planar graph $G$ and a ratio $r > 1$
  \State $k \gets \lceil 1/(r-1) \rceil$
  \State compute a planar embedding of $G$ and the level of every vertex
  \For{$i = 0, 1, \dots, k$}
    \State $V_i \gets$ all vertices whose level is \emph{not} congruent to $i$
           modulo $k+1$
    \State $G_i \gets G[V_i]$, which is $k$-outerplanar
    \State solve maximum independent set on $G_i$ exactly by
           \cref{lem:k-outerplanar-is}
  \EndFor
  \State \Return the largest of the $k+1$ independent sets found
\end{algorithmic}
\end{algorithm}

\begin{theorem}\label{thm:planar-is-ptas}
\Cref{alg:planar-is-ptas} runs in $\bigoh{8^k k n}$ time and has approximation
ratio $r$. Maximum independent set on planar graphs therefore has a PTAS.
\end{theorem}

\begin{proof}
Deleting the vertices of every $(k+1)$st level cuts the embedding into rings of
$k$ consecutive levels, so each component of $G_i$ is $k$-outerplanar and
\cref{lem:k-outerplanar-is} applies; there are $k+1$ iterations, giving the
running time. Since $k = \lceil 1/(r-1)\rceil$ is a constant once $r$ is fixed,
this is polynomial.

For the ratio, let $I$ be a maximum independent set of $G$. The $k+1$ deleted
classes $V \setminus V_i$ are pairwise disjoint, so they partition $I$ into
$k+1$ parts, and the smallest part has at most $|I|/(k+1)$ vertices. For that
$i$, the set $I \cap V_i$ is an independent set of $G_i$ of size at least
$|I| - |I|/(k+1) = |I| \cdot k/(k+1)$, and the exact solve on $G_i$ does at
least as well. Hence
\[ \frac{\OPT}{\mathrm{ALG}} \;\le\; \frac{|I|}{|I| \cdot k/(k+1)}
   \;=\; \frac{k+1}{k} \;=\; 1 + \frac1k \;\le\; r, \]
the last step because $k \ge 1/(r-1)$.
\end{proof}

Note what the algorithm does with the running time it is given: the exponent of
$8$ is $k \approx 1/(r-1)$, so pushing the ratio towards $1$ costs
exponentially. That is characteristic of a scheme that is a PTAS and not an
FPTAS --- and the next section shows the cost is unavoidable.

%==============================================================================
\section{The polynomial-bound technique}
\label{sec:polynomial-bound}
%==============================================================================

Where the gap technique separates the bottom of the scale, one observation
separates FPTAS from PTAS. It is almost too simple to be called a technique,
and it applies to a remarkable number of problems.

\begin{theorem}\label{thm:polynomial-bound}
Let $\Pi$ be an optimisation problem with integer solution values on which
$\OPT \le p(n)$ for a polynomial $p$, where $n$ is the input size. If $\Pi$ has
an FPTAS, then $\Pi$ can be solved exactly in polynomial time.
\end{theorem}

\begin{proof}
Run the FPTAS with $\varepsilon = 1/p(n)$. Its running time is polynomial in $n$
and in $1/\varepsilon = p(n)$, hence polynomial in $n$.

For a maximisation problem the guarantee gives
$\OPT/\mathrm{ALG} \le 1 + 1/p(n) = (p(n)+1)/p(n)$, so
\[ \mathrm{ALG} \;\ge\; \OPT \cdot \frac{p(n)}{p(n)+1}
   \;=\; \OPT - \frac{\OPT}{p(n)+1} \;>\; \OPT - 1, \]
using $\OPT \le p(n)$ in the last step. Since solution values are integers and
$\mathrm{ALG} \le \OPT$, this forces $\mathrm{ALG} = \OPT$. The minimisation
case is identical.
\end{proof}

\begin{corollary}\label{cor:no-fptas-is}
Unless $\mathrm{P} = \mathrm{NP}$, maximum independent set has no FPTAS --- not
even restricted to planar graphs, where it has a PTAS by
\cref{thm:planar-is-ptas}.
\end{corollary}

\begin{proof}
An independent set has at most $n$ vertices, so $\OPT \le n$ and
\cref{thm:polynomial-bound} applies with $p(n) = n$. Independent set is
\textrm{NP}-hard on planar graphs.
\end{proof}

So there is a genuine gap between levels~2 and~3 of the scale, and
\cref{thm:planar-is-ptas} sits strictly at level~3. The only property used is
that the optimum is polynomially bounded, which is why the argument is called
the \emph{polynomial-bound technique} and why it applies so widely: any problem
that asks for a subset of the input has $\OPT \le n$ for free. The problems that
escape it are the ones whose optima are numbers rather than counts --- like
\lang{Knapsack}, whose optimum can be exponentially large, which is exactly why
\cref{thm:knapsack-fptas} was possible.

%==============================================================================
\section{\texorpdfstring{$f(n)$}{f(n)}-APX: greedy set cover}
\label{sec:set-cover}
%==============================================================================

At the bottom of the scale the ratio is allowed to grow with the input. The
standard example is also one of the most useful algorithms in the subject.

\begin{probdef}{Set Cover}
  \instance A finite set $U$ with $|U| = n$ and a family $F$ of subsets of $U$
  whose union is $U$.
  \question A subfamily $C \subseteq F$ with $\bigcup_{S \in C} S = U$ and
  $|C|$ as small as possible.
\end{probdef}

\begin{algorithm}[ht]
\caption{Greedy \lang{Set Cover}}
\label{alg:greedy-set-cover}
\begin{algorithmic}[1]
  \State $C \gets \emptyset$
  \While{$U \ne \emptyset$}
    \State pick $S \in F$ maximising $|S \cap U|$
    \State $C \gets C \cup \{S\}$, \quad $U \gets U \setminus S$
  \EndWhile
  \State \Return $C$
\end{algorithmic}
\end{algorithm}

\begin{theorem}\label{thm:greedy-set-cover}
\Cref{alg:greedy-set-cover} is an $(\ln n + 1)$-approximation for
\lang{Set Cover}.
\end{theorem}

\begin{proof}
It runs in polynomial time and returns a cover, since the union of $F$ is $U$
and each round covers at least one new element.

Write $S_1, S_2, \dots$ for the sets picked in order, and let
$C_i = S_1 \cup \dots \cup S_{i-1}$ be the elements covered before $S_i$ is
picked. Charge each element $x \in U$ a cost: if $x$ is first covered by $S_i$,
put
\[ c(x) = \frac{1}{|S_i \setminus C_i|}. \]
Every set picked distributes a total charge of exactly $1$ among the elements it
newly covers, so
\[ \mathrm{ALG} = |C| = \sum_{x \in U} c(x). \]

Now bound the charge accumulated inside any single set $S \in F$, picked or not.
Write $S = \{s_1, \dots, s_m\}$ with the elements ordered so that $s_m$ is
covered first by the algorithm and $s_1$ last. Suppose $s_i$ is first covered at
the step that picks $S_j$. At that moment $s_1, \dots, s_i$ are all still
uncovered, so $|S \setminus C_j| \ge i$; and the algorithm picked $S_j$ because
it covers the most uncovered elements, so
$|S_j \setminus C_j| \ge |S \setminus C_j| \ge i$. Hence $c(s_i) \le 1/i$ and
\[ \sum_{x \in S} c(x) \;\le\; \sum_{i=1}^{m} \frac{1}{i} \;=\; H(|S|), \]
where $H(m) = \sum_{i=1}^m 1/i$ is the $m$th harmonic number, and
$H(m) \le \ln m + 1$ by comparison with $\int_1^m \mathrm{d}x/x$.

Finally let $C^{*}$ be an optimal cover. Every element lies in at least one set
of $C^{*}$, so
\[ \mathrm{ALG} = \sum_{x \in U} c(x)
   \;\le\; \sum_{S \in C^{*}} \sum_{x \in S} c(x)
   \;\le\; |C^{*}| \cdot H\bigl(\max_{S \in F} |S|\bigr)
   \;\le\; \OPT \cdot (\ln n + 1). \qedhere \]
\end{proof}

The technique --- charge the algorithm's cost to the elements, then re-collect
the charges along an optimal solution --- is worth remembering on its own. It is
the standard way of analysing a greedy algorithm when no single quantity $x$ in
the sense of \cref{sec:approx-proof-shape} presents itself.

This ratio cannot be improved: \citet{Feige1998} showed that approximating
\lang{Set Cover} within $(1-o(1))\ln n$ is impossible unless every problem in
\textrm{NP} has a slightly superpolynomial algorithm. So \lang{Set Cover} is at
level~5 of the scale and not at level~4.

%==============================================================================
\section{Independent set is not in APX}
\label{sec:is-not-apx}
%==============================================================================

One more separation, proved by a technique we have not used before: a
self-reduction, which takes an approximation algorithm and makes it better than
it was.

\begin{definition}[Graph product]\label{def:graph-product}
For a graph $G = (V,E)$, the graph $G \times G$ has vertex set $V \times V$,
with $(u,v)$ and $(u',v')$ adjacent if and only if either $u = u'$ and
$\{v,v'\} \in E$, or $\{u,u'\} \in E$.
\end{definition}

\begin{lemma}\label{lem:graph-product-is}
$G$ has an independent set of size $k$ if and only if $G \times G$ has one of
size $k^2$. Moreover, an independent set $J$ of $G \times G$ can be converted in
polynomial time into an independent set of $G$ of size at least $\sqrt{|J|}$.
\end{lemma}

\begin{proof}
If $I$ is independent in $G$ with $|I| = k$, then $I \times I$ is independent in
$G \times G$: two distinct elements $(u,v),(u',v')$ of it have either $u \ne u'$,
and then $\{u,u'\} \notin E$ while the first case of \cref{def:graph-product}
does not apply either, or $u = u'$ and $v \ne v'$ with $\{v,v'\} \notin E$.

Conversely let $J$ be independent in $G \times G$ and let
$U = \{u : (u,v) \in J \text{ for some } v\}$ be its projection. Then $U$ is
independent in $G$: if two distinct $u,u' \in U$ were adjacent, any
$(u,v),(u',v') \in J$ would be adjacent by the second case of
\cref{def:graph-product}. Likewise, for each fixed $u$ the fibre
$\{v : (u,v) \in J\}$ is independent in $G$, by the first case.

If $|U| \ge \sqrt{|J|}$ we return $U$. Otherwise some fibre contains at least
$|J|/|U| > \sqrt{|J|}$ vertices, and we return it. Taking $|J| = k^2$ gives an
independent set of size at least $k$ in $G$, which together with the first
paragraph proves the equivalence.
\end{proof}

\begin{lemma}[Self-reduction]\label{lem:is-self-reduction}
If maximum independent set has a polynomial-time $c$-approximation, then it has
a polynomial-time $\sqrt{c}$-approximation.
\end{lemma}

\begin{proof}
Given $G$, build $G \times G$ --- it has $n^2$ vertices, so this is polynomial
--- and run the $c$-approximation on it, obtaining an independent set $J$. By
\cref{lem:graph-product-is} the optimum of $G \times G$ is $\OPT^2$, where
$\OPT$ is the optimum of $G$, so $\OPT^2 / |J| \le c$. Convert $J$ into an
independent set of $G$ of size at least $\sqrt{|J|}$. Then
\[ \frac{\OPT}{\sqrt{|J|}} = \sqrt{\frac{\OPT^2}{|J|}} \le \sqrt{c}. \qedhere \]
\end{proof}

\begin{theorem}\label{thm:is-not-apx}
If maximum independent set is in \textrm{APX}, then it has a PTAS.
\end{theorem}

\begin{proof}
Let a $c$-approximation be given and let $\varepsilon > 0$. Applying
\cref{lem:is-self-reduction} $t$ times gives a $c^{1/2^t}$-approximation, and
$c^{1/2^t} \to 1$, so for some constant $t$ depending only on $c$ and
$\varepsilon$ we have $c^{1/2^t} \le 1+\varepsilon$. The construction is applied
a constant number of times, each squaring the number of vertices, so the final
graph has $n^{2^t}$ vertices --- polynomial for fixed $t$.
\end{proof}

\Cref{ch:hardness} proves that maximum independent set is \textrm{APX}-hard, and
hence has no PTAS unless $\mathrm{P} = \mathrm{NP}$; with
\cref{thm:is-not-apx} that says maximum independent set is not in \textrm{APX}
at all. It is in $f(n)$-APX for $f(n) = n/(\log n)^2$, and nothing much better
is possible: \citet{Zuckerman2007} showed that approximating it within
$n^{1-\delta}$ is \textrm{NP}-hard for every $\delta > 0$.

Things can be worse still. \Cref{thm:tsp-no-approx} rules out constant ratios
for \lang{TSP} without the triangle inequality, and the same construction with
$n\,p(n) + 1$ in place of $nc + 1$ on the non-edges rules out ratio $p(n)$ for
every polynomial $p$: that problem is outside $f(n)$-APX for every polynomial
$f$.

%==============================================================================
\section{Steiner tree}
\label{sec:steiner-tree}
%==============================================================================

The last section applies both halves of the chapter to one problem: the pool
exercise \emph{Approximating Steiner Tree}.

\begin{probdef}{Steiner Tree}
  \instance A connected undirected graph $G = (V,E)$ and a set $T \subseteq V$
  of \emph{terminals}; the vertices of $N = V \setminus T$ are available as
  intermediate vertices.
  \question A smallest set $S \subseteq E$ such that $(V(S),S)$ is a tree
  containing every terminal, where $V(S)$ denotes the set of endpoints of $S$.
\end{probdef}

With $T = V$ this is the minimum spanning tree problem, solvable exactly; with
$|T| = 2$ it is a shortest path. In between it is \textrm{NP}-hard, and the
reason is that one must decide which non-terminals to use.

\subsection{A 2-approximation}

The idea is to pretend the non-terminals do not exist, by hiding them inside
edge costs.

\begin{definition}[Metric closure]\label{def:metric-closure}
Let $G_T = (T,F)$ be the complete graph on the terminals, where the cost
$c(\{t_1,t_2\})$ of an edge is the length of a shortest $t_1$--$t_2$ path in
$G$, and let $P(\{t_1,t_2\}) \subseteq E$ be the edges of such a path.
\end{definition}

\begin{algorithm}[ht]
\caption{MST heuristic for \lang{Steiner Tree}}
\label{alg:steiner-mst}
\begin{algorithmic}[1]
  \State build the metric closure $G_T$ and record a shortest path $P(f)$ for
         every edge $f \in F$
  \State compute a minimum spanning tree $M \subseteq F$ of $G_T$ with respect
         to $c$
  \State \Return $\bigcup_{m \in M} P(m)$
\end{algorithmic}
\end{algorithm}

\begin{theorem}\label{thm:steiner-2approx}
\Cref{alg:steiner-mst} runs in polynomial time and is a $2$-approximation for
\lang{Steiner Tree}.
\end{theorem}

\begin{proof}
All-pairs shortest paths and a minimum spanning tree are polynomial, so the
algorithm is.

\emph{It returns a feasible solution.} The union of the paths connects every
pair of terminals that $M$ connects, and $M$ spans $T$, so all terminals lie in
one connected component of the returned edge set. If that set contains a cycle,
delete edges until it is a tree; this only decreases the cost.

\emph{It costs at most $2\OPT$.} Let $S^{*}$ be an optimal Steiner tree, of size
$\OPT$. Walk around $S^{*}$ in a depth-first traversal, using every edge once in
each direction: a closed walk of length exactly $2\OPT$ visiting all terminals.
List the terminals in the order the walk first meets them, $t_1, \dots, t_r$,
and consider the cycle $t_1 t_2 \cdots t_r t_1$ in $G_T$. Each of its edges
costs at most the length of the corresponding piece of the walk, since $c$ is
the shortest path distance, and the pieces are disjoint parts of the walk, so
the cycle costs at most $2\OPT$. Deleting its most expensive edge leaves a
spanning tree of $G_T$ of cost at most $2\OPT$, so $c(M) \le 2\OPT$. Finally the
returned edge set has size at most $\sum_{m \in M} |P(m)| = c(M)$, since paths
may share edges but never add any.
\end{proof}

\begin{intuition}
This is the double-tree argument of \cref{thm:tsp-double-tree} run backwards. In
metric TSP we converted a tree into a tour by doubling; here we convert an
optimal tree into a tour in order to get a cheap spanning tree of the closure.
Both directions cost the same factor $2$, and both pay it in the same place:
traversing every edge twice.
\end{intuition}

\subsection{No FPTAS}

The second half of the pool question asks for the limits, and the previous
sections supply them almost immediately.

\begin{theorem}\label{thm:steiner-no-fptas}
Unless $\mathrm{P} = \mathrm{NP}$, \lang{Steiner Tree} has no FPTAS.
\end{theorem}

\begin{proof}
A Steiner tree is the edge set of a tree on at most $|V|$ vertices, so
$\OPT \le |V| - 1 \le n$, and solution values are integers.
\Cref{thm:polynomial-bound} with $p(n) = n$ turns an FPTAS into an exact
polynomial-time algorithm, and \lang{Steiner Tree} is \textrm{NP}-hard.
\end{proof}

The weaker statement --- no approximation within an absolute constant --- has
its own short proof, in the style of \cref{thm:knapsack-no-additive}, and is
worth giving because it does not need \cref{thm:polynomial-bound}.

\begin{proposition}\label{prop:steiner-no-additive}
Unless $\mathrm{P} = \mathrm{NP}$, no polynomial-time algorithm approximates
\lang{Steiner Tree} within an absolute difference $c$, for any integer constant
$c \ge 1$.
\end{proposition}

\begin{proof}
Given an instance $(G,T)$, subdivide every edge of $G$ into a path of $c+1$
edges, the new internal vertices being non-terminals. Call the result $G'$.

Every Steiner tree of $G'$ may be assumed to use each subdivided path either
entirely or not at all: an internal vertex of a path has degree $2$ in $G'$ and
is not a terminal, so a partially used path ends in a vertex of degree $1$ that
is not a terminal, and deleting that edge keeps the solution feasible and makes
it smaller. Hence Steiner trees of $G'$ correspond to Steiner trees of $G$ with
all sizes multiplied by $c+1$, and $\OPT(G') = (c+1)\OPT(G)$.

All attainable values in $G'$ are multiples of $c+1$, so a solution within $c$
of optimal is optimal, and the assumed algorithm would solve the
\textrm{NP}-hard problem \lang{Steiner Tree} exactly in polynomial time.
\end{proof}

%==============================================================================
\section{Notes and further reading}
\label{sec:approx-ii-further-reading}
%==============================================================================

This chapter follows Johan van Rooij's lecture slides. The classification of
approximation qualities and the accompanying negative techniques are the
subject of \citet[Chapter~3]{Ausiello1999}; \citet{WilliamsonShmoys2011} and
\citet{Vazirani2001} are the standard texts on the algorithms.

The FPTAS for \lang{Knapsack} is \citet{IbarraKim1975}. The PTAS for planar
graphs is \citet{Baker1994}, and the treewidth bound behind
\cref{lem:k-outerplanar-is} is \citet{Bodlaender1988}. Greedy set cover and its
$\ln n$ analysis were found independently by \citet{Johnson1974},
\citet{Lovasz1975} and \citet{Chvatal1979}, the last in the weighted setting;
the matching lower bound is \citet{Feige1998}. The graph-product
self-reduction for independent set is classical and is given in
\citet{GareyJohnson1979}; the best possible ratio is \citet{Zuckerman2007}. \Cref{ex:planar-colouring} rests on the
four colour theorem, \citet{AppelHaken1977}, and \cref{ex:vizing} on
\citet{Vizing1964}.

%==============================================================================
\section{Exercises}
\label{sec:ex-approx-ii}
%==============================================================================

\begin{exercise}
Show that every FPTAS is a PTAS, and that every problem with a PTAS is in
\textrm{APX}. Give a problem separating PTAS from \textrm{APX}, quoting whatever
hardness result you need.
\end{exercise}

\begin{exercise}
The proof of \cref{thm:no-additive-is} takes $c+1$ disjoint copies. Carry the
same argument out for \lang{Vertex Cover}, and say which step needs changing.
\end{exercise}

\begin{exercise}
In \cref{alg:knapsack-fptas} the values are scaled and the weights are not.
What goes wrong if you scale the weights instead?
\end{exercise}

\begin{exercise}
\Cref{thm:knapsack-fptas} assumes $\varepsilon \le 1$. What ratio does the same
choice of $c$ give for $\varepsilon > 1$, and why does the restriction cost
nothing?
\end{exercise}

\begin{exercise}
Give an FPTAS for \lang{Subset Sum}: given positive integers $a_1,\dots,a_n$ and
a target $t$, find a subset with sum at most $t$ and as large as possible.
\end{exercise}

\begin{exercise}
Show that the analysis of \cref{alg:greedy-set-cover} is essentially tight:
construct instances on which the greedy algorithm returns a cover of size
$\Omega(\log n) \cdot \OPT$.
\end{exercise}

\begin{exercise}
Adapt \cref{alg:planar-is-ptas} to \lang{Vertex Cover} on planar graphs. Which
direction of the ratio argument changes, and why is \emph{adding} the deleted
level class the right move for a minimisation problem?
\end{exercise}

\begin{exercise}
\Cref{lem:graph-product-is} says that the independence number satisfies
$\alpha(G \times G) = \alpha(G)^2$. Show that for the \emph{disjoint union}
$\alpha(G_1 \cup G_2) = \alpha(G_1) + \alpha(G_2)$, and explain why
\cref{lem:is-self-reduction} needs the product and \cref{thm:no-additive-is}
needs the union.
\end{exercise}

\begin{exercise}
In \cref{prop:steiner-no-additive} the edges are subdivided into paths of length
$c+1$. Why can the same effect not be achieved by simply declaring each edge to
have weight $c+1$, in the problem as stated here?
\end{exercise}
